\documentclass[Total.tex]{subfiles}
\begin{document}
	\chapter{Representation of Compact Groups}
		\section{Basic Representation Theory}
			\subsection{Definitions}
			\subsubsection{The Topology of $E$}
			\subsubsection{Representation of Topological Group}
			\begin{Def}
				Let $G$ be a topological group, and $E$ be a \underline{topological vector space}.
				\begin{itemize}
					\item We say that: $E$ is a $G$-module, if there is an action
					\begin{gather*}
						G\times E\rightarrow E\qquad (g,x)\mapsto gx
					\end{gather*}
					\item 
					an action $(g,x)\mapsto gx:G\times E\rightarrow E$ such that:
					\begin{itemize}
						\item $:x\mapsto gx,E\rightarrow E$ is a \textbf{continuous vector space endomorphism} of $E$, for each $g\in G$;
						\item $g\mapsto gx,G\rightarrow E$ is continuous for each $x\in E$;
					\end{itemize}
					\item A \textbf{representation} of $G$ is a continuous map
					\begin{gather*}
						\pi:G\rightarrow \mathcal{L}_p(E)
					\end{gather*}
					
					satysfying
					\begin{gather*}
						\pi(1)=id_E\qquad \pi(gh)=\pi(g)\pi(h);g,h\in G
					\end{gather*}
				\end{itemize}
			\end{Def}
			
			\textbf{Remark} \begin{itemize}
				\item  Possible Representation
				\begin{gather*}
					\pi:G\rightarrow \mathcal{L}_p(E)\qquad g\mapsto (x\mapsto gx)
				\end{gather*}
				\item If $\pi:G\rightarrow \mathcal{L}_p(E)$ is a representation, then the function
				\begin{gather*}
					G\times E\rightarrow E\qquad (g,x)\mapsto gx:=\pi(g)(x)
				\end{gather*}
				
				endows $E$ with the structure of a $G$-module.
			\end{itemize}
			
			We could have another famous statement:  It is the theory of $G$-module.
			\begin{theorem}
				Assume that: $E$ is a $G$-module for a topological group $G$, and that
				\begin{gather*}
					\pi:G\rightarrow \mathcal{L}_p(E)
				\end{gather*}
				
				is the associated representation. 
				
				If: $E$ is a Baire space, then for every compact subspace $K$ of $G$, the set
				\begin{gather*}
					\pi(K)\subseteq Hom(E,E)
				\end{gather*}
				
				is equi-continuous at $0$. That is: For any neighborhood $V$ of $0$ in $E$, there is a neighborhood $U$ of $0$ such that
				\begin{gather*}
					KU\subseteq V
				\end{gather*}
				
				As a consequence, if $G$ is locally compact, the function
				\begin{gather*}
					G\times E\rightarrow E\qquad (g,x)\mapsto gx
				\end{gather*}
				
				is continuous.
			\end{theorem}
			\begin{proof}
				内容...
			\end{proof}
			
			\begin{example}
				内容...
			\end{example}
			
			\begin{corollary}
				内容...
			\end{corollary}
			\begin{proof}
				内容...
			\end{proof}
			\subsection{Haar Integral}
			
			Take the following notion: Given a measure $\mu$,
			\begin{gather*}
				\mu(f):=\int_X fd\mu=(\mu,f)
			\end{gather*}
			
			or indeed $=\int_G f(g)d\mu(g)$. A special case is the classical Fourier analysis
			\begin{gather*}
				\mathbb{T}=\RR/\ZZ
			\end{gather*}
			
			\subsubsection{Invariant Measure}
			
			Take the following notation $(_gf)(x):=f(g^{-1}x);g\in G$(multiplicative group).
			
			\begin{Def}
				\begin{itemize}
					\item Let $G$ denote a compact group.
					
					A measure $\mu$ is called \textbf{invariant} if
					\begin{gather*}
						\mu(_g f)=\mu(f),\forall g\in G;f\in E=C(G,\mathbb{K})
					\end{gather*}
					\item A measure $\mu$ is called \textbf{Haar measure}, if it is invariant and positive.
					
					Positive: $\mu(f)\geq 0,\forall f\geq 0$.
					\item The measure is \textbf{normalized} if $\mu(1)=1$.
				\end{itemize}
			\end{Def}
			
			\begin{example}
				Consider the morphism
				\begin{gather*}
					p:\mathbb{R}\rightarrow \mathbb{T}\qquad t\mapsto t+\mathbb{Z}
				\end{gather*}
			\end{example}
			
			\begin{example}
				内容...
			\end{example}
			\subsubsection{Measure on Group}
				Preparations for existence theorem:
				\begin{lemma}
					内容...
				\end{lemma}
				\begin{proof}
					内容...
				\end{proof}
				
				
				\begin{theorem}
					(The existence and Uniqueness Theorem of Haar Measure) For each \textbf{compact group}
				\end{theorem}
				\begin{proof}
					内容...
				\end{proof}
				
				\begin{corollary}
					内容...
				\end{corollary}
				\begin{proof}
					内容...
				\end{proof}
			\subsection{Consequence of Haar Measure}
				\begin{theorem}
					(Weyl's Trick) 
					\begin{itemize}
						\item Let $G$ be a compact group, $E$ a $G$-module which is also a Hilbert space.(not Hilbert $G$-module) 
						
						Then, there is a scalar product relative to which all operators $\pi(g)$ are unitary.
						\item Specifically, if $(-,-)$ is the given scalar product on $E$, then
						\begin{gather*}
							<x|y>=\int_G (gx|gy)dg
						\end{gather*}
						
						defines a scalar product such that
						\begin{gather*}
							M^{-2}(x|x)\leq \left\langle x|x\right\rangle \leq M^2(x|x)\qquad M=sup\{\sqrt{(gx|gx)}|g\in G,(x|x)\leq 1\}
						\end{gather*}
						
						and that 
						\begin{gather*}
							\langle gx| gy\rangle=\langle x| y\rangle;x,y\in E,g\in G
						\end{gather*}
					\end{itemize}
				\end{theorem}
				\begin{proof}
					Firstly, the new-defined product is well-defined. Then, since Haar measure is positive, we have $(gx|gx)\geq 0\Rightarrow \langle x|x\rangle\geq 0$. Thus, the constant $M$ is well-defined.
					
					Then,
					\begin{gather*}
						内容...
					\end{gather*}
					
					Finally:
					\begin{gather*}
						<hx|hy>=\int_G (ghx|ghy)dg=\int_G (gx|gy)dg=<x,y>\qquad\mbox{(Invariance of Haar Measure)}
					\end{gather*}
				\end{proof}
				\subsubsection{Hilbert $G$-module}
				\begin{Def}
					(Hilbert $G$-module) If $G$ is a \underline{topological space}. Then, a \textbf{Hilbert $G$-module} is a \underline{Hilbert space} $E$ and a $G$-module such that: All operators $\pi(g)$ is unitary. i.e.
					\begin{gather*}
						(gx|gy)=(x,y);x,y\in E,g\in G
					\end{gather*}
				\end{Def}
				
				\begin{example}
					Let $G$ be a compact group and $\mathcal{H}_0$ the vector space $C(G,\mathcal{K})$ equipped with the scalar product
					\begin{gather*}
						(f_1|f_2)=\gamma(f_1\bar{f}_2)=\int_G f_1(g)\overline{f_2(g)}dg
					\end{gather*}
					
					Indeed:
					\begin{itemize}
						\item 
						\item 
					\end{itemize}
					
					Then, the open set $U=\{t\in G|(f\bar{f})(t)>0\}$ contains $g$. Hence is nonempty.
				\end{example}
				\subsection{The Main Theorem on Hilbert Modules for Compact Groups}
					Recall:
					\begin{itemize}
						\item A \textbf{sesquilinear form} is a function
						\begin{gather*}
							B:\mathcal{H}\times \mathcal{H}\rightarrow K
						\end{gather*}
						that is linear in the first and conjugate linear in the second. (in the sense of $\CC$-linear)
						\item  It is \textbf{bounded} if 
						\begin{gather*}
							|B(x,y)|\leq M\Vert x\Vert\cdot \Vert y\Vert\qquad x,y\in \mathcal{H}
						\end{gather*}
						\item A way to define sesquilinear form: Suppose $T$ be a \textbf{bounded linear operator} on $\mathcal{H}$. Then,
						\begin{gather*}
							B(x,y)=(Tx|y)
						\end{gather*}
						
						defines a sesquilineaar from with $M=\Vert T\Vert$: In the view of Cauchy-Schwartz
						\begin{gather*}
							|(x,y)|\leq \Vert x\Vert\cdot \Vert y\Vert 
						\end{gather*}
					\end{itemize}
					But conversely:
					\begin{lemma}
						If $B$ is a sesquilinear form, then there exists a unique bounded operator $T$ of $\mathcal{H}$ such that
						\begin{gather*}
							\Vert T\Vert\leq M\qquad B(x,y)=(Tx|y)
						\end{gather*}
					\end{lemma}
					\begin{proof}
						Such $T$ is specified by:
						
						Uniqueness:
					\end{proof}
					
					Under this trick, we could define another operator, which is stabil under group actions:
					\begin{lemma}
						Let $G$ denote a compact grooup and $T$ a bounded operator on a Hilbert $G$-module $E$. Then: There exists a unique bounded operator $\tilde{T}$ on $E$, with $\Vert\tilde{T}\Vert\leq \Vert T\Vert$ such that
						\begin{gather*}
							(\tilde{T}x|y)=\int_G (Tgx|gy)dg=\int_G  (\pi(g)^{-1}T\pi(g)(x)|y)dg
						\end{gather*}
					\end{lemma}
					\begin{proof}
						内容...
					\end{proof}
					
					\begin{proposition}
						(Criterion for $\tilde{S}$) The following statements are equivalent for an operator $S$ of $H$:
						\begin{itemize}
							\item $S\in Hom_G(H,H)$;
							\item $S=\tilde{S}$;
							\item There is an operator $T$ such that $S=\tilde{T}$;
						\end{itemize}
					\end{proposition}
					\begin{proof}
						内容...
					\end{proof}
					\subsubsection{$G$-Submodules}
					\begin{Def}
						If $G$ is a topological group, $E$ a $G$-module, then a vector-subspace $V$ of $E$ is called a \textbf{submodule} if
						\begin{gather*}
							GV\subseteq V
						\end{gather*}
						
						also called \textbf{invariant subspace}.
					\end{Def}
					\begin{lemma}
						内容...
					\end{lemma}
					\begin{proof}
						内容...
					\end{proof}
					
					\begin{lemma}
						内容...
					\end{lemma}
					\begin{proof}
						内容...
					\end{proof}
				\subsection{The Fundamental Theorem on Unitary Modules}
				\subsection{Definition of a Compact Lie Group,Compact Group by Compact Lie Groups}
				
				Given in Lie Theory.
				
		\section{Theorem of Peter and Weyl}
		\section{G-Module Theory}
			\subsection{Vector Valued Integration}
		\section{Compact Groups-Hopf Algebras}
	\chapter{Characters}
	\chapter{Lie Group}
		These will be given in the note of Lie group and Algebras;
		\begin{itemize}
			\item Linear Lie Groups;
			\item Compact Lie Groups;
		\end{itemize}
	\chapter{Abelian Groups}
\end{document}